\begin{proposition}[P4.25]\hypertarget{tag-02WX}{}\label{night:P4.25}\label{stmt:P4.25}
\StatementRecord{night:P4.25}{proved}{02WX}{\SeriesRef{night:P4.16}, \SeriesRef{night:P4.20}, \SeriesRef{night:P4.6}, \SeriesRef{night:P4.8}}

Every semi-log-canonical reduced fibre support in
Part~4.2 has normal crossings. In particular, let
$\rho:T\to N\to S$ be a resolution of the normalization
of an irreducible support, and let $D$ be the integral
divisor defined by
$\omega_T(D)\simeq\rho^*\omega_S$. When this actual
divisor is reduced and has normal crossings, $S$
has normal crossings, its normalization is smooth,
and P4.6--P4.8 apply.
\end{proposition}
\begin{proof}

\begin{enumerate}[label={$\langle1\rangle$\arabic*.},leftmargin=*]
\item We first explain why the stated resolved-conductor
case is semi-log-canonical. At a generic conductor
point, an exponent one on every normalization
branch forces the transverse plane curve to be an
ordinary node. Indeed the conductor exponent on
branch $i$ is
$2\delta_i+\sum_{j\ne i}I_{ij}$; positivity of
intersection multiplicities leaves exactly two
smooth transverse branches. The support $S$ is a
Cartier hypersurface, hence is Cohen--Macaulay and
has invertible dualizing sheaf. On its normal
normalization $N$, the pair consisting of $N$
and its reduced conductor is log canonical.
To check this last assertion, pull its canonical
divisor to $T$. The resulting boundary is exactly
$D$, by the displayed isomorphism and the
divisorial orders of a local generator. A smooth
surface with reduced normal-crossing boundary is
log canonical: blowing up a point gives new
boundary coefficient $\mu-1\le1$, and this
bound persists under successive point blowups.
Discrepancy comparison therefore gives the
assertion on $N$. This is precisely the
semi-log-canonical definition of
Liu--Rollenske, Definition~2.3 and Lemma~2.2,
which explicitly applies to complex analytic
spaces.

\item We now treat all local semi-log-canonical
hypersurface germs. We use the complete analytic
classification in Liu--Rollenske,
\emph{Two-dimensional semi-log-canonical
hypersurfaces}, main Theorem, Table~1,
Propositions~3.1 and~4.1. The normal singular
germs are isolated. Any such germ is excluded
in our family: a small generic linear
perturbation of its reduced equation has
Morse critical points, the number counted
with multiplicity being the positive
colength of its gradient ideal. The equality
of this number before and after perturbation
follows by conservation of the local
intersection multiplicity of the three
gradient equations. A constant perturbation
puts their critical values away from zero.
The original zero divisor is smooth near
a small transverse boundary sphere, so it
retains that boundary property under these
perturbations. The proper boundary
transport proved in P4.20 carries a
vanishing sphere of square $-2$ into a
nearby torus fibre, which is impossible.

\item For the nonnormal germs we give the exact
deformations and check their boundary strata.
The ordinary double point and $xyz=0$
already have normal crossings. A pinch point
$x^2+y^2z=0$ is excluded by P4.16.
The remaining double-point equations are
\[
 x^2+y^2z^2,\qquad
 x^2+y^2(z^2+y^{q-2}),\quad q\ge3.
\]
For the first equation the two normalization
branches are the smooth graphs $x=\pm i yz$.
Their tangent planes coincide at the origin,
whereas they meet transversely off that point
along the two axes. The local proof of P4.20
therefore supplies a vanishing sphere after
a deformation retaining its boundary
ordinary-double strata.

\item For finite $q$ put
\[
 G_s=x^2+y^2(z^2+y^{q-2}+s).
\]
For $s\ne0$ small, at each point
$(0,0,z_0)$ with $z_0^2=-s$ the function
$Z=z^2+y^{q-2}+s$ is a local coordinate:
its $z$ derivative is $2z_0\ne0$.
The actual equation there is exactly
$x^2+y^2Z$, a pinch point. At the original
boundary double axis $x=y=0$, $z\ne0$,
the transverse Hessian determinant is
$4(z^2+s)$, which stays nonzero on the
compact boundary when $s$ is small.
All boundary smooth points remain smooth.
Thus the zero divisor retains the same
boundary smooth and double strata.
The spherical class of square $-4$
from the pinch-point nearby fibre transports,
by the detailed P4.20 construction, to
the original nearby torus fibre.
This excludes every finite $q$.

\item The remaining triple-point equations are
\[
 xyz+x^p,\quad p\ge3,\qquad
 xyz+x^p+y^q,\quad q\ge p\ge3.
\]
For the first put
$G_s=xyz+x^p+s x^2
=x(yz+x^{p-1}+s x)$.
When $s\ne0$, its two local normalization
branches are smooth: the derivative in $x$
of the second factor at the origin is $s$.
Their tangent planes coincide, so the
origin is not a normal crossing. Along
the boundary $y$ axis the transverse
Hessian determinant is $-y^2$, and along
the boundary $z$ axis it is $-z^2$.
Both axes persist as ordinary-double
strata. The local argument of P4.20 on
the deformed germ supplies a compact
spherical class of nonzero square, and
the same boundary transport gives the
contradiction.

\item For the second equation put
\[
 G_s=xyz+x^p+y^q+s x^2+s y^3,
 \qquad 0<|s|\ll1.
\]
Its transverse Hessian along the boundary
$z$ axis has determinant $-z^2$,
so its sole boundary double stratum
persists. We check the smooth normalization
of its germ at the origin explicitly.
Since $\partial^2G_s/\partial x^2(0)=2s$,
the critical equation in $x$ has a unique
analytic solution $x=\xi(y,z)$ with
$\xi$ divisible by $yz$. The parametric
holomorphic Morse change in $x$ writes
\[
 G_s=X^2+y^2H(y,z).
\]
Divisibility follows directly on substituting
$\xi$: each of $\xi yz,\xi^p,s\xi^2$
is divisible by $y^2$, as are the two
pure $y$ terms. Moreover
\[
 H_y(0,0)=
 \begin{cases}1+s&q=3,\\s&q>3,\end{cases}
 \qquad H(0,z)=-z^2/(4s).
\]
The finite
birational map
\[
 \{v^2+H(y,z)=0\}\longrightarrow\{G_s=0\},
 \qquad (v,y,z)\longmapsto(X=yv,y,z),
\]
followed by the inverse Morse coordinate
change, is the normalization. Its source
is smooth because $H_y(0,0)\ne0$,
and the map is finite because $v$
satisfies the displayed monic equation.
It is birational on $y\ne0$ and is
therefore the normalization by normality
of its source. At the origin its
differential has rank one, so this germ
is not a normal crossing. P4.20 again
produces a spherical class, with the
already verified boundary transport.

\item These calculations exhaust the analytic
classification. Thus only the smooth,
ordinary-double, and triple-normal-crossing
germs remain. Their local equations give
the claimed normal-crossing support and
smooth normalization.

\end{enumerate}
\end{proof}